%Desenvolvido por: Prof. MsC. Arthur Monteiro Filho 
%Versão:1.0.0
%Data07/05/2023
%Professor do IFES Campus São Mateus 
%e-mail:arthur.monteiro@ifes.edu.br


\documentclass[aspectratio=169]{beamer}
\usepackage[brazil]{babel}
\usepackage[utf8]{inputenc}
\usepackage[T1]{fontenc}
\usepackage{textcomp}
\usepackage{xcolor}
%\usepackage[alf]{abntex2cite}	% Citações padrão ABNT
\usepackage{color}			    % Controle das cores
\usepackage{graphicx}			% Inclusão de gráficos
\usepackage{txfonts}			% Fontes virtuais
\usepackage{hyperref}
\usepackage{multimedia}
\setbeamercovered{transparent}

% ---

\title{Título da Aula}
\date[]{\today}
\author[]{Nome do Professor}

\usetheme{texsx}

\begin{document}
	
	\begin{frame}
		\titlepage
		\thispagestyle{empty}	%não numerar o slide CAPA	
	\end{frame}
%%%%%%%%%%%%%%%%%%%%%%%%%%%%%%%%%%%%%%%%%%%%%%%%%%%%%%%%%%%%%%%%%%%%%%%%%%%%%%%%%%%%%%%%%%%%%	
	\section[Sumário]{}	
	\begin{frame}
		\frametitle{Sumário}
		\tableofcontents
	\end{frame}
%%%%%%%%%%%%%%%%%%%%%%%%%%%%%%%%%%%%%%%%%%%%%%%%%%%%%%%%%%%%%%%%%%%%%%%%%%%%%%%%%%%%%%%%%%%%%
	\section{Objetivos}
	
	\begin{frame}{Objetivos}
		\begin{itemize}
			\item Um
				\begin{itemize}		
					\item a
					\item aa		
				\end{itemize}
			\item Dois
		\end{itemize}
	\end{frame}
%%%%%%%%%%%%%%%%%%%%%%%%%%%%%%%%%%%%%%%%%%%%%%%%%%%%%%%%%%%%%%%%%%%%%%%%%%%%%%%%%%%%%%%%%%%%%	
	\section{Introdução}
	
	\begin{frame} 
		\frametitle{Introdução}					
	\end{frame}
%%%%%%%%%%%%%%%%%%%%%%%%%%%%%%%%%%%%%%%%%%%%%%%%%%%%%%%%%%%%%%%%%%%%%%%%%%%%%%%%%%%%%%%%%%%%%
	\section{Conclusão}
	 \begin{frame}{Conclusão}
		\begin{itemize}
		\uncover<1>{\item one}
		\uncover<2>{\item two}
		
			
		\end{itemize}
	\end{frame}
%%%%%%%%%%%%%%%%%%%%%%%%%%%%%%%%%%%%%%%%%%%%%%%%%%%%%%%%%%%%%%%%%%%%%%%%%%%%%%%%%%%%%%%%%%%%%	
	\section{Recomendação de Leitura}
	
	\begin{frame}
		\frametitle{Leitura Recomendada}
		
		
		\begin{itemize}
			\item Capítulo 1 - da página 1 a 10;
			
		\end{itemize}
		
		% TODO: \usepackage{graphicx} required
		\begin{figure}
			\centering
			%\hspace{8cm} \vspace{0cm}
			\includegraphics[width=0.28\linewidth]{imagens/CAPA-livro-geracao-termeletrica}
			\caption{Fonte:Geração Termelétrica\cite{lora2004geraccao}}
			\label{fig:capa-livro-geracao-termeletrica}
		\end{figure}
		
	\end{frame}
%%%%%%%%%%%%%%%%%%%%%%%%%%%%%%%%%%%%%%%%%%%%%%%%%%%%%%%%%%%%%%%%%%%%%%%%%%%%%%%%%%%%%%%%%%%%
	\section{Próxima Aula}

	\begin{frame}
		\frametitle {Próxima Aula}
		\begin{itemize}			
			\item Na próxima aula estudaremos...			
		\end{itemize}
		
    

%		\movie{\includegraphics[width=0.28\linewidth]{imagens/sala-de-aula.png}}{imagens/sala-de-aula.mp4}
	\end{frame}
	
%%%%%%%%%%%%%%%%%%%%%%%%%%%%%%%%%%%%%%%%%%%%%%%%%%%%%%%%%%%%%%%%%%%%%%%%%%%%%%%%%%%%%%%%%
\section{Referências}

% --- O comando \allowframebreaks ---
% Se o conteúdo não se encaixa em um quadro, a opção allowframebreaks instrui 
% beamer para quebrá-lo automaticamente entre dois ou mais quadros,
% mantendo o frametitle do primeiro quadro (dado como argumento) e acrescentando 
% um número romano ou algo parecido na continuação.

\begin{frame}[allowframebreaks]{Referências}
%	\bibliographystyle{abnt-alf}
%	\bibliography{abntex2-modelo-references}
  \bibliographystyle{apalike}
  \bibliography{bibliografia}
\end{frame}	


%%%%%%%%%%%%%%%%%%%%%%%%%%%%%%%%%%%%%%%%%%%%%%%%%%%%%%%%%%%%%%%%%%%%%%%%%%%%%%%%%%%%%%
\setbeamertemplate{navigation symbols}{}
\usebackgroundtemplate{\includegraphics[width=\paperwidth]{imagens/SLIDE2.png}}

\begin{frame}	 
		
	\begin{center}
	\uncover<1> \Huge{OBRIGADO!}\\
	\end{center} 	
 
\centering
	 \begin{center}
	 	\textbf{E-MAIL}: 
	 	 %\href{mailto:}{}\\descomentar e adicionar o seu e-mail.
	 \textbf{Horário de Atendimento}:\\
	 Quinta-feira: 13h00 às 14:00.
	 \end{center}
	 	 
\end{frame}	
%%%%%%%%%%%%%%%%%%%%%%%%%%%%%%%%%%%%%%%%%%%%%%%%%%%%%%%%%%%%%%%%%%%%%%%%%%%%%%%%%%%%%%
%último slide, não mexer!	
%Acrescenta o slide do final da apresentação institucional 	
%	{%<--- Start local changes
		\setbeamertemplate{navigation symbols}{}
		\usebackgroundtemplate{\includegraphics[width=\paperwidth]{imagens/ultCAPA.png}}
		\begin{frame}[plain]
		\end{frame}
%	}%<---- Finish local changes

%%%%%%%%%%%%%%%%%%%%%%%%%%%%%%%%%%%%%%%%%%%%%%%%%%%%%%%%%%%%%%%%%%%%%%%%%%%%%%%%%%%%%% 
	
\end{document}